本章汇总对原文全部命题、仿真数值与实证数值所做的独立核验，并列出原文推导与表述中需要修改的问题。核验结果写成三类判定。通过表示符号恒等或数值落在预设容差内。偏差表示结论方向成立，但数值、条件或措辞与原文表述有出入。不通过表示按预设判据不成立。问题按严重度分三级。高表示结论或论证不成立，答辩或审稿时会被直接追问。中表示结论成立，但条件、近似或数字需要修正。低表示编号、符号、措辞或排版问题。

\subsection{核验的范围与方法}

核验共 199 项，覆盖模型设定、命题 1--6、定义 1、推论 1、更新过程仿真与全部实证数值，每个命题至少有一项。所用方法分三种。符号推导用计算机代数系统逐式验证原文的代数结论，例如负二项抽稀的概率母函数、Fisher 信息、全方差分解与 AR(1) 双重求和的闭式。数值复算按原文给出的参数重算原文报告的每个数字，例如交叉点 $m_\times=53.08$、经验平衡点 $b/a=22.6471$、偏差项 0.2669 与各类份额。蒙特卡洛模拟在原文的生成机制下重复抽样，比较模拟矩与解析公式，每项不超过 $2\times10^5$ 次抽样，随机数流固定，全部可复现。

总体结果为通过 187 项、偏差 10 项、不通过 2 项。表~\ref{tab:ck-overview} 按主题列出覆盖范围与结果。

\begin{table}[!tbp]
\centering
\small
\caption{核验项按主题的覆盖范围与判定结果}
\label{tab:ck-overview}
\setlength{\tabcolsep}{4pt}
\begin{tabular}{L{2.8cm}L{5.4cm}L{2.4cm}L{0.9cm}L{0.9cm}L{0.9cm}L{1.0cm}}
\toprule
主题 & 核验内容与原文位置 & 方法 & 条数 & 通过 & 偏差 & 不通过 \\
\midrule
命题 1（符号） & 负二项抽稀封闭性、得分函数、Fisher 信息与 CRLB（原文式 1--6） & 符号推导 & 13 & 13 & 0 & 0 \\
命题 1（数值） & 抽稀分布、估计量方差与 Fisher 信息 & 数值复算、蒙特卡洛 & 15 & 15 & 0 & 0 \\
命题 2、定义 1、命题 3 & 方差分解 $v_R+A/m$、交叉点与翻倍降幅、杠杆拆分（原文式 7--13） & 符号推导、蒙特卡洛 & 22 & 21 & 1 & 0 \\
命题 4 & 预算权衡、AR(1) 有效期数、可行期数集合（原文式 14--17） & 符号推导、蒙特卡洛 & 18 & 16 & 2 & 0 \\
推论 1 & 成本最优单期规模、角点解、整数解与 AR(1) 修正（原文式 18） & 符号推导、整数枚举 & 18 & 17 & 1 & 0 \\
命题 5 & 真实感染与报告病例的单步预测方差分解（原文式 21--22） & 符号推导、蒙特卡洛 & 11 & 10 & 0 & 1 \\
命题 6 & 单调似然比、精确检验水平、功效与 Wald 水平（原文式 23--25） & 符号推导、精确计算 & 24 & 24 & 0 & 0 \\
更新过程仿真 & 相对误差恒等式、条件偏差–方差分解、原文图 3 数值（原文式 27--29） & 符号推导、蒙特卡洛 & 19 & 19 & 0 & 0 \\
实证算术一致性 & 仿射拟合、份额分解、阶段偏差、外部验证（原文式 30--31 与原文表 1--6） & 数值复算 & 36 & 32 & 3 & 1 \\
随机分母与报告过离散 & 分母噪声对 $a$、$b$、$m_\times$ 的影响（原文式 31） & 蒙特卡洛、delta 法 & 14 & 11 & 3 & 0 \\
FGLS 与聚类 Bootstrap & 一致性、比值偏差、区间覆盖率 & 蒙特卡洛 & 9 & 9 & 0 & 0 \\
\midrule
合计 & & & 199 & 187 & 10 & 2 \\
\bottomrule
\end{tabular}
\end{table}

按主题看，核验结论可以概括为三条。第一，命题 1--6 与推论 1 的全部代数结论经符号推导成立。负二项抽稀后的分布与 $\mathrm{NB}(nk,p^*)$ 的总变差距离在 $10^{-15}$ 量级。估计量方差的模拟值与原文式 6 之比为 1.0199、1.0039、0.9830，得分函数方差与 Fisher 信息之比与之同步。$m=53$ 时环境项份额为 0.499，与定义 1 的 50\% 一致。Wald 检验在 $m=5$、$k=0.1$ 时的精确水平为 0.0781，原文的蒙特卡洛值 0.081 落在其误差带内。蒙特卡洛与解析公式在 $|z|<3$ 内一致，唯一越界项见下一小节，属抽样波动。第二，需要改写的是近似与表述：$T_{\mathrm{eff}}$ 近似恒高估环境项；精确检验的水平是“不超过 $\alpha$”；原文图 2 报告的“首次”达到 80\% 功效的样本量之后功效会跌回，稳定样本量多 1.4\%--3.0\%。按初始规模排列时 $v_R+A/m$ 不是下界。第三，实证映射存在系统偏离：截距与斜率分别是情景隐含值的 3.5--6.9 倍与 3.2--3.7 倍，二项分母噪声约解释斜率放大的一半，截距放大随平均回溯延迟单调增加。FGLS 在均值设定正确时无偏，州聚类 Bootstrap 无明显欠覆盖（仿真规模受限）。

\subsection{未通过与偏差项}

表~\ref{tab:ck-nonpass} 逐条列出 12 项非通过结果。其中多数来自讲义补充的核验（原文没有给出对应数字），判定为“偏差”表示原文结论方向成立、但需要补充条件或修正措辞。

\begin{table}[!tbp]
\centering
\footnotesize
\caption{12 项非通过核验结果的原文值、复核值与解释}
\label{tab:ck-nonpass}
\setlength{\tabcolsep}{3pt}
\begin{tabular}{L{2.5cm}L{1.6cm}L{2.3cm}L{2.7cm}L{0.9cm}L{4.6cm}}
\toprule
主题 & 原文位置 & 原文值 & 复核值 & 判定 & 解释 \\
\midrule
住院情景下杠杆非对称的数量级 & 第 467、613、658 行 & 作为实证启示引用，未给数量 & $A/\bar R\in[1.019,1.180]$ & 偏差 & 报告率约 1\% 时两条杠杆近乎等效，引用命题 3 时应说明（I-05） \\
$T_{\mathrm{eff}}$ 近似的方向 & 第 254--256 行 & 未说明方向 & 所有 $(T,r)$ 下均高估环境项 & 偏差 & 高估比例为 $2r(1-r^T)/[T(1-r)(1+r)-2r(1-r^T)]$（I-07） \\
$T\gg1/(1-r)$ 条件下的误差 & 第 252--256 行 & 可良好近似 & $T=26$、$r=0.8$ 时高估 20.5\% & 偏差 & 满足 $T>1/(1-r)=5$ 不足以保证良好近似，宜给误差界或精确闭式（I-07） \\
情景交叉点范围 & 第 515 行 & 21.3--49.3 & 21.28--49.18 & 偏差 & 49.3 来自 $A$ 中取 $v_R=0.08$、分母取 0.04 的混用，差异很小但口径应统一（I-08） \\
命题 5 的 $\operatorname{Var}(C)$，$\mathbb{E}[\Lambda]=40$，$\mathrm{CV}(\Lambda)=0.3$ & 原文式 22 & 理论 37.49 & 模拟 36.97，$z=-4.45$ & 不通过 & 蒙特卡洛波动；同配置的 $\operatorname{Var}(I)$ 通过，20 组独立重复的相对误差为 $-0.07\%\pm0.27\%$ \\
误差份额的归一化口径 & 第 551 行对第 554 行 & 54.8\%/45.2\% & 去残差重新归一可复现；第 554 行为 54.6\%/45.1\% & 偏差 & 同一规格用了两种归一化，正文未注明（I-18） \\
覆盖率骤降的增长率离散 & 第 613 行 & 扩张 192.1\%，即 2.92 倍 & 均匀漏报上界 1.29 倍 & 偏差 & 远超上界，支持结构性扰动解释，与命题 3 只能定性联系（I-21） \\
截距与情景隐含值 & 第 508 行 & 0.0991 与 0.014--0.029 相符 & 为隐含值的 3.5--6.9 倍 & 不通过 & 原文表述错误，公式本身无误（I-14） \\
斜率与情景隐含值 & 第 508、515 行 & 未比较；隐含 0.610--0.707 & 2.2434，为 3.2--3.7 倍 & 偏差 & 斜率同样被放大，$m_\times$ 的量级吻合来自分子分母同向放大（I-14） \\
分母含噪时 $b/a$ 与 $m_\times$ & 第 512、515 行 & 以 $b/a$ 对照情景 $m_\times$ & 模拟 $b/a=59.2$，真值 $A/v_R=30.7$ & 偏差 & 分母噪声只放大斜率，$b/a$ 高估 $m_\times$（I-13、I-14） \\
二项分母噪声能解释的 $b$ & 第 506--512 行 & 未讨论，$b=2.2434$ 直接作规模项 & 预测 1.1551（实证的 51\%）；无分母噪声 0.6601（29\%） & 偏差 & 剩余缺口可由 $\varphi\approx1.96$ 的报告过离散补足，属可能解释之一（I-13） \\
报告过离散 $\varphi=1.96$ 下的 $b$ & 第 506--512 行 & 未讨论 & 预测 2.2386；模拟 2.0190（SE 0.0343），$a=0.0227$ & 偏差 & delta 法在噪声大时高估约 5\%--10\%；截距停在 $v_R/\bar R^2$，其放大需另找来源（I-12） \\
\bottomrule
\end{tabular}
\end{table}

两项“不通过”的性质完全不同。第一项是命题 5 的报告病例方差在 $\mathbb{E}[\Lambda]=40$、$\mathrm{CV}(\Lambda)=0.3$ 下的蒙特卡洛比较，模拟值 36.97 对理论值 37.49，$z=-4.45$。这是抽样波动，原文的公式没有错误。同一配置下真实感染方差 $\operatorname{Var}(I)$ 通过。用 20 组独立随机数流各 $10^4$ 次重复复核，$\operatorname{Var}(C)$ 的相对误差平均为 $-0.07\%$，波动幅度 $0.27\%$，与零无差异。单组模拟中 $\Lambda$ 的重尾使批次方差偏大，$|z|$ 偶尔越过 3。第二项是原文第 508 行的措辞错误：实证截距 0.0991 被写成与情景隐含值 0.014--0.029“相符”，实际为 3.5--6.9 倍。这一项涉及的公式与数值本身都能复现，错的是比较结论，修改方法见下一小节的 I-14。

\subsection{原文推导与表述问题}

逐行对照原文后共整理出 24 个问题：高严重度 4 个（I-12、I-13、I-14、I-23），中严重度 12 个，低严重度 8 个。表~\ref{tab:ck-issues-a}、表~\ref{tab:ck-issues-b} 与表~\ref{tab:ck-issues-c} 依次列出 I-01 至 I-24，行号均为原文 \LaTeX{} 源码的行号。

\begin{table}[!tbp]
\centering
\footnotesize
\caption{原文推导与表述问题（I-01 至 I-08）}
\label{tab:ck-issues-a}
\setlength{\tabcolsep}{3pt}
\begin{tabular}{L{0.9cm}L{0.8cm}L{1.9cm}L{5.6cm}L{5.8cm}}
\toprule
编号 & 严重度 & 原文行号 & 问题 & 修改建议 \\
\midrule
I-01 & 低 & 756、761（对照 310、336） & 附录“命题 5/6”两个小节的标题与正文编号互换；正文交叉引用指向正确，错的是标题文字。 & 交换两个小节的标题，答辩材料统一按正文编号。 \\
I-02 & 中 & 89--121、688 & 命题 1 依赖的完整负二项分布与独立性假设没有集中写出，正文只给前两阶矩，附录直接使用全分布。 & 在第 94 行后增加“假设 1”：条件负二项分布、同期共享 $R_t$、$k$ 与 $\rho$ 已知、报告与传播条件独立。 \\
I-03 & 中 & 145、196、419、715 & CRLB 在每个有限 $n$ 上精确达到，原文写作“收敛”；UMVUE 缺少论证；两者都要求 $k$ 已知。 & 第 419 行改为“比值 0.988--1.017 在蒙特卡洛误差内，与 CRLB 精确达到一致”；附录补 Lehmann--Scheffé 论证。 \\
I-04 & 中 & 111、114、244 & $\mathbb{E}[R_t\mid n]=\bar R$ 在时间序列中不自动成立：$n_t$ 由过去传播决定，环境自相关时与 $R_t$ 相关。 & 第 114 行后说明该假设对设计固定的 $n$ 成立，对实证中的 $D_t$ 需另行讨论，并在局限性中呼应。 \\
I-05 & 中 & 212--231、467、613、658 & 命题 3 的非对称性在住院情景下只有 1.02--1.18 倍，数量上很弱。 & 命题 3 后给出 $A/\bar R=1+\rho(\bar R^2+v_R)/(k\bar R)$，说明住院监测中两条杠杆近似等价。 \\
I-06 & 中 & 168 对 617 & “无法随横截面聚合消除”与全国汇总降低环境方差的结论矛盾；$v_R$ 只在单位共享同一 $R_t$ 时不可约。 & 第 168 行改为“无法通过同一环境内扩大规模消除，跨相关性不完全的环境聚合可削减其 $(1-\bar c)$ 部分”。 \\
I-07 & 中 & 254--256、753 & $T_{\mathrm{eff}}$ 近似恒高估环境项，$T=26$、$r=0.8$ 时高估 20.5\%，$T=10$ 时高估 65.8\%。 & 正文给出附录第 751 行的精确闭式；或把条件改为 $T\gg2r/(1-r^2)$ 并写明近似偏保守。 \\
I-08 & 中 & 273--283 & 推论 1 的数值来源未写；上端 49.3 应为 49.18；连续松弛与整数解、AR(1) 近似未讨论。 & 写明取第 515 行情景，49.3 改为 49.2；补“$T$ 须为整数，预算较小时应枚举”。 \\
\bottomrule
\end{tabular}
\end{table}

\begin{table}[!tbp]
\centering
\footnotesize
\caption{原文推导与表述问题（I-09 至 I-16）}
\label{tab:ck-issues-b}
\setlength{\tabcolsep}{3pt}
\begin{tabular}{L{0.9cm}L{0.8cm}L{1.9cm}L{5.6cm}L{5.8cm}}
\toprule
编号 & 严重度 & 原文行号 & 问题 & 修改建议 \\
\midrule
I-09 & 低 & 299--312 & 命题 5 的两式是任意分布下的全方差恒等式，没有给出可计算的结论；第 312 行孤句与上下文无衔接。 & 增加负二项分支、$\Lambda\perp R$ 下的推论 $v_R+A/m+\mathrm{CV}^2(\Lambda\mid\mathcal F)(\bar R^2+v_R)$；删除或改写第 312 行。 \\
I-10 & 中 & 323、340、358、364 & 精确检验的实际水平应为“不超过 $\alpha$”；非随机化检验只在它的实际水平下一致最优；用 $\hat k$ 时水平无精确保证；“首次”达到 80\% 功效的样本量之后功效会跌回。 & 第 340 行改“不超过”；注明随机化检验 \citep{karlin1956theory}；局限性中讨论 $\hat k$；报告稳定样本量 2107/581/235/585/1684 与精确 Wald 水平。 \\
I-11 & 低 & 375；图 1(c) & $\max(1,I)$ 截断改变了生成过程，触发频率未报告；图 1(c) 标题渲染为“m imes”。 & 报告各规模组的截断次数；绘图标签改用原始字符串。 \\
I-12 & 高 & 441--446、506--512 & 周增长比 $G$ 跨越约 2 个以上世代，与单代 $R_t$ 口径不同；截距随平均滞后单调增大（0.0682/0.0991/0.1524）。 & 第 446 行后说明口径差，把 $a$ 解读为“周增长比的规模无关离散”，并报告三种权重下的 $a$。 \\
I-13 & 高 & 506--512 & $D_t$ 由报告病例构造，二项噪声约解释斜率的一半，并定量复现三种权重下 $b$ 的比值（1.44/0.86 对 1.40/0.86）。 & 在局限性中给出 delta 法公式与三种权重的对照，写明 $b$ 不能直接读作 $A/\bar R^2$。 \\
I-14 & 高 & 508、515 & $a$、$b$ 分别是情景隐含值的 3.5--6.9 倍与 3.2--3.7 倍，“相符”不成立；$b/a$ 的吻合来自同向放大。 & 删去“相符”，写明倍数与来源；第 515 行说明 $b/a$ 的吻合不构成结构识别。 \\
I-15 & 低 & 251、312 & 符号 $A$ 同时指“目标 A”与规模相关方差系数。 & 估计目标改称“目标 I/目标 II”，或方差系数改记 $\kappa$。 \\
I-16 & 低 & 254、330--332、418 & 符号 $r$ 同时指 AR(1) 系数与负二项形状参数 $nk$。 & 形状参数统一写 $nk$，或 AR(1) 系数改记 $\phi$。 \\
\bottomrule
\end{tabular}
\end{table}

\begin{table}[!tbp]
\centering
\footnotesize
\caption{原文推导与表述问题（I-17 至 I-24）}
\label{tab:ck-issues-c}
\setlength{\tabcolsep}{3pt}
\begin{tabular}{L{0.9cm}L{0.8cm}L{1.9cm}L{5.6cm}L{5.8cm}}
\toprule
编号 & 严重度 & 原文行号 & 问题 & 修改建议 \\
\midrule
I-17 & 低 & 多处（如 176、291、409） & 若干展示式有 PDF 式号但没有 label，正文无法引用，只能写“上式”。 & 被引用的式子补 label，不被引用的改为无编号环境。 \\
I-18 & 低 & 551 对 554 & 误差份额 54.8\%/45.2\% 与 54.6\%/45.1\% 的归一化口径不同，正文未注明。 & 第 551 行注明“按两项之和归一”，或统一用第 554 行口径；可补份额比 $=m_\times^{\mathrm{emp}}/D_{\mathrm{harm}}$，$D_{\mathrm{harm}}\approx27.4$。 \\
I-19 & 中 & 515；中期报告 & 工具变量的排除限制无法检验；工具变量不消除被解释变量中的分母噪声；“排除内生性”的说法过强。 & 沿用原文“替代测度下的稳健性”定位，中期报告与答辩口径去掉“排除内生性”。 \\
I-20 & 中 & 617 & 全国汇总只有 $N=28$，样本平均 0.0348 而中位数 0.0115，分布偏斜；隐含 $\bar c\approx0.34$ 未报告。 & 报告反推的 $\bar c$ 与 Bootstrap 区间，说明样本量限制。 \\
I-21 & 中 & 613 & 增长率离散扩张 2.92 倍，远超均匀漏报的上界 1.29 倍。 & 加入该上界作为定量论据，支持结构性扰动判断；与命题 3 的联系改为定性表述。 \\
I-22 & 低 & 187、196、409 等 & “拐点”用词不当：$V''=2A/m^3>0$，$V$ 严格凸，没有拐点。 & 统一为“方差占比交叉点”或“平衡点”，给出弹性 $-m_\times/(m+m_\times)$。 \\
I-23 & 高 & 393、406 & “始终不低于单代下界”没有被证明；按初始规模排列时比值为 0.75--0.98，反而低于参照。 & 改为“按实现规模分箱时本仿真中不低于参照曲线（1.01--1.17 倍）”，并说明这不是一般性下界。 \\
I-24 & 中 & 508、512 & 用 CV² 随规模递减论证 $1/\hat\mu^2$ 权重，逻辑方向反了。 & 改写为“CV² 随规模下降说明该权重未达最优；均值设定正确时权重只影响效率”，补 OLS 对照或改进权重。 \\
\bottomrule
\end{tabular}
\end{table}

\subsubsection{I-12：周增长比与世代再生数的口径差}

理论中的 $R_t$ 是一代传播的再生数，命题 2 的环境方差 $v_R$ 是单代量。原文实证的增长代理 $G_{j,t+1}=C_{j,t+1}/\allowbreak D_{j,t}$ 是周度比值，分母又是滞后 1--3 周的加权和，主权重下平均回溯 1.45 周。流感世代间隔约 3 天，一周约跨 2 个世代，所以 $G$ 衡量的是跨越多个世代的累积增长，$\ln G\approx\sum_g\ln R_g$。各州各周增长趋势的异质性随跨度累积，且不随规模衰减，于是进入截距。决定性的证据来自原文表 1：平均滞后 1.0、1.45、2.0 周时，截距依次为 0.0682、0.0991、0.1524，单调增大；而在理想映射下截距等于 $v_R/\bar R^2$，与权重无关，蒙特卡洛也证实二项分母噪声不改变截距。原文第 508 行把截距与 $v_R/\bar R^2$ 直接比较，忽略了这一口径差。修改时宜在第 446 行后写明 $G$ 与 $R_t$ 的口径差，把 $a$ 解读为“周增长比的规模无关离散”，并把三种权重下的 $a$ 作为证据报告。详细推导见第~\ref{sec:em-intercept} 小节。

\subsubsection{I-13：随机分母对斜率的放大}

原文式 30 的 $D_{j,t}$ 由报告病例构造，本身带有二项报告噪声，违背理想映射中“分母无噪”的条件。由一阶 delta 法，斜率近似为 $b\approx[A+(1-\rho)(\bar R^2+v_R)\sum w^2]/\bar R^2$，第二项正比于滞后权重的平方和。三种权重的 $\sum w^2$ 为 0.495、1、1/3，预测的斜率比值为 1.44 与 0.86，原文表 1 的实证比值为 1.40 与 0.86，蒙特卡洛为 1.311 与 0.836。水平上，二项分母噪声把斜率从 0.6601（实证的 29\%）抬到 1.1551（实证的 51\%），再加 $\varphi\approx1.96$ 的报告过离散可补足剩余缺口。分母噪声只抬高斜率，所以 $b/a$ 会高估 $m_\times$：模拟中 59.2 对真值 30.7。这一机制把原文表 1 中经验交叉点对滞后权重的敏感性变成了可解释的结果，可以写成论文的新增结论。修改时宜在局限性中给出上式与三种权重下的定量对照，并写明 $b$ 不能直接读作 $A/\bar R^2$。详细推导见第~\ref{sec:em-delta} 小节。

\subsubsection{I-14：第 508 行“相符”不成立}

原文第 515 行的流感住院情景（$\bar R\approx1.67$，$k\in[0.2,0.5]$，$\rho\in[0.0056,0.021]$，$v_R\in[0.04,0.08]$）隐含截距 $v_R/\bar R^2\in[0.0143,0.0287]$，隐含斜率 $A/\bar R^2\in[0.610,0.707]$。实证截距 0.0991 是隐含值的 3.5--6.9 倍，实证斜率 2.2434 是隐含值的 3.2--3.7 倍，原文第 508 行“经验值 0.0991 与参数空间隐含值 0.014--0.029 相符”不成立，这是核验中仅有的两项“不通过”之一。$b/a=22.6$ 落在情景区间 21.3--49.3 内，原因是截距与斜率被放大了相近的倍数，而放大机制各不相同：截距的放大来自 I-12 的口径差，斜率的放大来自 I-13 的分母噪声，倍数接近只是巧合。所以第 515 行“经验平衡点在数量级上与之相符”也不能当作结构验证。建议第 508 行改为“截距高于情景隐含值（3.5--6.9 倍），差额来自 $G$ 与世代 $R$ 的口径差及代理偏差”，第 515 行改为“$b/a$ 与情景区间同量级，但 $a$、$b$ 各自均高于隐含值，这一吻合不构成结构识别”。答辩时若被问到经验交叉点与理论的关系，宜直接给出两个倍数，再说明两种放大机制。

\subsubsection{I-23：“始终不低于单代下界”没有被证明}

原文第 393 行与第 406 行称更新过程的经验误差“始终不低于单代静态下界”，依据是更新过程的条件偏差–方差分解，原文给出的分解如下式所示：
\begin{equation*}
\mathbb{E}\left[(G_{t+1} - \bar{R})^2 \;\middle|\; D_t > 0\right] = \mathbb{E}\left[\frac{\sigma_{t+1 \mid t}^2}{D_t^2} + \left(\frac{\mu_{t+1 \mid t}}{D_t} - \bar{R}\right)^2 \;\middle|\; D_t > 0\right].
\tag{原文式 28}
\end{equation*}
该式本身经符号推导与模拟核验成立（$m=6$ 至 200 两端全部一致）。问题出在由它推出“下界”的一步。右端第二项非负，只能推出 $\mathbb{E}[(G-\bar R)^2]\ge\mathbb{E}[\sigma^2/D^2]$；要得到 $\mathbb{E}[(G-\bar R)^2]\ge v_R+A/m$，还需要 $\mathbb{E}[\sigma^2/D^2]\ge v_R+A/m$。在负二项分支下，条件方差除以 $(\rho\,\mathbb{E}[\Lambda])^2$ 才等于 $v_R+A/m$ 加一个非负项，而原文式 28 除的是报告历史 $D^2$；增长期中 $D$ 与 $\rho\,\mathbb{E}[\Lambda]$ 的比值是随机的，不等式方向无法确定。数值上，结论取决于横轴口径：按初始生成参数 $m$ 排列时，经验与参照之比为 0.75--0.98（核验复算为 0.77--0.92），低于 1；按实际实现规模分箱时为 1.01--1.17（核验复算为 1.03--1.12）。评价周的实际规模中位数为初始值的 1.7--2.4 倍（复算 1.82--2.36 倍），这解释了两种口径的差别。因此“不低于下界”只能算实现规模口径下的经验观察，尚无定理支撑；把偏差项非负当作代数证明的论证不成立。建议改为“按实现规模分箱时，本仿真中经验误差均不低于参照曲线（比值 1.01--1.17）；这与原文式 28 中非负偏差项的方向一致，但这不构成一般性下界”。
